% ============================================================
%  ACADEMIC CV TEMPLATE
%  Suitable for: PhD students, researchers, academics
%
%  LaTeX concepts covered:
%    \documentclass, \usepackage, geometry, fontenc, inputenc,
%    \section*, \textbf, \textit, \hfill, \hrule, \vspace,
%    itemize (nested), hyperref, \href, \url, \noindent
%
%  Instructions:
%    1. Replace all placeholder text with your own information.
%    2. Delete any sections you do not need.
%    3. Click Recompile (or press Ctrl+Enter) to update the PDF.
% ============================================================

% ------ DOCUMENT CLASS & PACKAGES ------
% \documentclass sets the overall type and base font size.
\documentclass[11pt, a4paper]{article}

% geometry lets you control the page margins precisely.
\usepackage[top=2cm, bottom=2cm, left=2.2cm, right=2.2cm]{geometry}

% fontenc and inputenc ensure accented characters print correctly.
\usepackage[T1]{fontenc}
\usepackage[utf8]{inputenc}

% hyperref turns URLs and email addresses into clickable links in the PDF.
% hidelinks removes the coloured boxes around links.
\usepackage[hidelinks]{hyperref}

% Remove paragraph indentation and add a small gap between paragraphs.
\setlength{\parindent}{0pt}
\setlength{\parskip}{4pt}

% ------ HELPER COMMAND ------
% This defines a reusable command for section headings with a horizontal rule.
% \newcommand{\cvsection}[1] means "one argument, referred to as #1".
\newcommand{\cvsection}[1]{%
  \vspace{10pt}%
  {\large \textbf{#1}}%        % #1 is replaced by whatever you pass in
  \vspace{3pt}%
  \hrule                        % draws a full-width horizontal rule
  \vspace{6pt}%
}

% ============================================================
%  DOCUMENT BEGINS HERE
% ============================================================
\begin{document}

% ------ HEADER ------
% \begin{center} ... \end{center} centres everything inside.
\begin{center}
  {\LARGE \textbf{Your Full Name}} \\[4pt]
  % \\ starts a new line. [4pt] adds 4 points of extra space below.
  Department of Your Department \\
  University Name, City, Country \\[6pt]

  % \href{url}{display text} makes a clickable link (from hyperref package).
  \href{mailto:your.email@university.edu}{your.email@university.edu}
  \quad $\cdot$ \quad               % \quad is a medium space; $\cdot$ is a bullet in math mode
  \href{https://yourwebsite.com}{yourwebsite.com}
  \quad $\cdot$ \quad
  \href{https://github.com/yourusername}{github.com/yourusername}
\end{center}

\vspace{4pt}


% ------ RESEARCH INTERESTS ------
\cvsection{Research Interests}

Write 1--2 sentences describing your main research area and questions.
For example: I'm broadly interested in supporting the safe and responsible use of AI by analyzing the norms, tensions, and subtle technical assumptions that underlie AI system development.


% ------ EDUCATION ------
\cvsection{Education}

% \textbf makes text bold. \textit makes it italic.
% \hfill pushes the date all the way to the right margin.
\textbf{Doctor of Philosophy in Your Field} \hfill \textit{Expected Month Year} \\{}
University Name, City, Country \\
Thesis: \textit{``Your Thesis Title''} \\
Supervisor: Prof.\ Supervisor Name

\vspace{8pt}

\textbf{Master of Science in Your Field} \hfill \textit{Month Year} \\{}
University Name, City, Country \\
Thesis: \textit{``Thesis Title''} \quad GPA: X.X / 4.0

\vspace{8pt}

\textbf{Bachelor of Science in Your Field} \hfill \textit{Month Year} \\{}
University Name, City, Country \quad GPA: X.X / 4.0


% ------ RESEARCH EXPERIENCE ------
\cvsection{Research Experience}

\textbf{Graduate Research Assistant} \hfill \textit{Month Year -- Present} \\
\textit{Lab / Group Name}, University Name \\
Supervisor: Prof.\ Name

% itemize creates a bullet list. Each \item is one bullet point.
% Note: when bullet text starts with [, wrap it in braces: \item {[text]}
% otherwise LaTeX reads [...] as a custom bullet label.

\begin{itemize}
  \setlength{\itemsep}{2pt}   % reduces space between items
  \item Describe your main research contribution or project in one sentence.
  \item Describe a method, dataset, or tool you developed or used.
  \item Mention any collaboration, publication outcome, or result.
\end{itemize}

\vspace{8pt}

\textbf{Research Intern} \hfill \textit{Month Year -- Month Year} \\
\textit{Organisation / Lab}, City, Country

\begin{itemize}
  \setlength{\itemsep}{2pt}
  \item Describe the project you worked on.
  \item Describe your contribution and any results.
\end{itemize}


% ------ PUBLICATIONS ------
\cvsection{Publications}

% enumerate creates a numbered list.
\begin{enumerate}
  \setlength{\itemsep}{4pt}

  \item \textbf{Your Name}, Co-Author Name, and Co-Author Name.
        ``\textit{Title of Your Paper}.''
        In \textit{Conference or Journal Name}, Year.
        \href{https://doi.org/xxxxx}{doi:xxxxx}

  \item \textbf{Your Name} and Co-Author Name.
        ``\textit{Title of Second Paper}.''
        In \textit{Journal Name}, Vol.\ X, No.\ Y, Year.

\end{enumerate}


% ------ TEACHING ------
\cvsection{Teaching}

\textbf{Teaching Assistant} --- Course Name (Course Code) \hfill \textit{Year} \\{}
University Name

\begin{itemize}
  \setlength{\itemsep}{2pt}
  \item Led weekly tutorial sessions of N students.
  \item Graded assignments and provided written feedback.
\end{itemize}


% ------ AWARDS \& HONOURS ------
\cvsection{Awards \& Honours}

% You can also use itemize for a simple list without descriptions.
\begin{itemize}
  \setlength{\itemsep}{2pt}
  \item \textbf{Award Name} \hfill \textit{Year}
  \item \textbf{Scholarship Name}, Awarding Body \hfill \textit{Year}
  \item \textbf{Fellowship Name} \hfill \textit{Year}
\end{itemize}


% ------ SKILLS ------
\cvsection{Skills}

% Nested itemize: an itemize inside another itemize creates a sub-list.
\begin{itemize}
  \setlength{\itemsep}{3pt}

  \item \textbf{Programming:}
    \begin{itemize}
      \item Proficient: Python, R, MATLAB
      \item Familiar: C++, JavaScript
    \end{itemize}

  \item \textbf{Tools \& Frameworks:} PyTorch, TensorFlow, Git, Docker, \LaTeX

  \item \textbf{Languages:} English (fluent), Other Language (native)

\end{itemize}


% ------ REFERENCES ------
\cvsection{References}

Available upon request.

% ============================================================
%  DOCUMENT ENDS HERE
% ============================================================
\end{document}
